\subsection{Ranks of projective modules}

DEFINITION 1. Let P be a finitely generated projective A-module. For every prime ideal p of A, the rank of the free \(A_{p}\) -module \(P_{p}\) is called the rank of P at p and is denoted by \(\operatorname{rg}_{\mathfrak{p}}(\mathbf{P})\) .

By Theorem 1 the integer-valued function \(\mathfrak{p} \mapsto \mathrm{rg}_{\mathfrak{p}}(\mathrm{P})\) is locally constant on \(\mathbf{X} = \operatorname{Spec}(\mathbf{A})\) ; it is therefore constant if \(\mathbf{X}\) is connected and in particular if the ring \(\mathbf{A}\) is an integral domain (§ 4, no. 3, Corollary 2 to Proposition 15).

DEFINITION 2. Let \(n\) be an integer \(\geqslant 0\) . A projective A-module \(P\) is said to be of rank \(n\) if it is finitely generated and \(\operatorname{rg}_{\mathfrak{p}}(P) = n\) for every prime ideal \(\mathfrak{p}\) of \(A\) .

Clearly every finitely generated free A-module L is of rank n in the sense of Definition 2, n being equal to the dimension (or rank) of L defined in Algebra, Chapter II, §7, no. 2.

A projective module of rank 0 is zero (§ 3, no. 3, Corollary 2 to Theorem 1). If A is not reduced to 0 and a projective A-module P is of rank n, the integer n is determined uniquely; it is then denoted by rg(P).

THEOREM 2. Let P be an A-module and n an integer \(\geqslant 0\) . The following properties are equivalent:

\begin{enumerate}
\def\labelenumi{(\alph{enumi})}
\item
  P is projective of rank n.
\item
  \(\mathbf{P}\) is finitely generated and, for every maximal ideal \(\mathfrak{m}\) of \(\mathbf{A}\) , the \(\mathbf{A}_{\mathfrak{m}}\) -module \(\mathbf{P}_{\mathfrak{m}}\) is free of rank \(n\) .
\item
  \(\mathbf{P}\) is finitely generated and, for every prime ideal \(\mathfrak{p}\) of \(\mathbf{A}\) , the \(\mathbf{A}_{\mathfrak{p}}\) -module \(\mathbf{P}_{\mathfrak{p}}\) is free of rank \(n\) .
\item
  For every maximal ideal m of A, there exists \(f \in A - m\) such that the \(A_{f}\) -module \(P_{f}\) is free of rank n.
\end{enumerate}

By Definition 2 and Theorem 1, (a) and (c) are equivalent; (b) implies (c), as, for every prime ideal p of A, there exists a maximal ideal m containing p and, writing \(p' = p_{m}\) , \(P_{p}\) is isomorphic to \((\mathrm{P}_{\mathfrak{m}})_{\mathfrak{p'}}\) (§ 2, no. 5, Proposition 11); if \(P_{m}\) is free of rank n, so then is \(P_{p}\) . Property (c) implies (d) by virtue of Theorem 1 and the fact that, if \(f \in A - m\) and \(m' = m_{f}\) , \(P_{m}\) is isomorphic to \((\mathrm{P}_{f})_{\mathfrak{m'}}\) and therefore the ranks of \(P_{f}\) and \(P_{m}\) are equal. Finally, this last argument and Theorem 1 show that (d) implies (b).

Remark. If A is an integral domain, a projective A-module admits a well-defined rank (in the sense of Definition 2), as has been observed above; moreover, this rank coincides with the rank defined in Algebra, Chapter II, §7, no. 2; it is sufficient to apply Theorem 2 (c) with p = (0).

Let E and F be two finitely generated projective A-modules. We know (Algebra, Chapter II, §§ 2 and 3) that \(\mathbf{E} \times \mathbf{F}\) , \(\mathbf{E} \otimes_{\mathbf{A}} \mathbf{F}\) , \(\operatorname{Hom}_{\mathbf{A}}(\mathbf{E}, \mathbf{F})\) and the dual \(\mathbf{E}^*\)

of E are projective and finitely generated; so is the exterior power \(\wedge^k\) E for every integer \(k > 0\) (Algebra, Chapter III). Also, it follows immediately from Definition 1 and §2, no. 7, Propositions 18 and 19 and no. 8, that, for every prime ideal \(\mathfrak{p}\) of A:

\[
\mathrm{rg} _ {\mathfrak {p}} (\mathrm{E} \times \mathrm{F}) = \mathrm{rg} _ {\mathfrak {p}} (\mathrm{E}) + \mathrm{rg} _ {\mathfrak {p}} (\mathrm{F})\tag{1}
\]

\[
\operatorname{rg} _ {\mathfrak {p}} (\mathrm{E} \otimes_ {\mathrm{A}} \mathrm{F}) = \operatorname{rg} _ {\mathfrak {p}} (\mathrm{E}). \operatorname{rg} _ {\mathfrak {p}} (\mathrm{F})\tag{2}
\]

\[
\operatorname{rg} _ {\mathfrak {p}} (\operatorname{Hom} _ {\mathbf {A}} (\mathrm{E}, \mathrm{F})) = \operatorname{rg} _ {\mathfrak {p}} (\mathrm{E}). \operatorname{rg} _ {\mathfrak {p}} (\mathrm{F})\tag{3}
\]

\[
\operatorname{rg} _ {\mathfrak {p}} (\mathrm{E} ^ {*}) = \operatorname{rg} _ {\mathfrak {p}} (\mathrm{E})\tag{4}
\]

\[
\operatorname{rg} _ {\mathfrak {p}} (\overset {k} {\wedge} \mathrm{E}) = \binom{\operatorname{rg} _ {\mathfrak {p}} (\mathrm{E})}{k}.\tag{5}
\]

If the ranks of E and F are defined, so are those of \(E \times F\) , \(E \otimes_{A} F\) , \(\operatorname{Hom}_{A}(E, F)\) , \(E^{*}\) and \(\Lambda^{k} E\) and the above equations also hold with the index p omitted. Moreover:

COROLLARY. For a finitely generated projective A-module P to be of rank n, it is necessary and sufficient that \(\wedge^{n}P\) be of rank 1.

PROPOSITION 4. Let B be a commutative A-algebra and P a projective A-module of rank n.~The B-module \(\mathbf{P}_{(\mathbf{B})} = \mathbf{B} \otimes_{\mathbf{A}} \mathbf{P}\) is then projective of rank n.

We know that \(P_{(B)}\) is projective and finitely generated (Algebra, Chapter II, §5, no. 1, Corollary to Proposition 4). If q is a prime ideal of B and p its inverse image in A, then

\[
\left(\mathrm{P} _ {(\mathrm{B})}\right) _ {\mathrm{q}} = \left(\mathrm{P} \otimes_ {\mathrm{A}} \mathrm{B}\right) \otimes_ {\mathrm{B}} \mathrm{B} _ {\mathrm{q}} = \mathrm{P} \otimes_ {\mathrm{A}} \mathrm{B} _ {\mathrm{q}} = \left(\mathrm{P} \otimes_ {\mathrm{A}} \mathrm{A} _ {\mathrm{p}}\right) \otimes_ {\mathrm{A}} \mathrm{B} _ {\mathrm{q}}
\]

and, as, by hypothesis, \(\mathbf{P} \otimes_{\mathbf{A}} \mathbf{A}_{\mathfrak{p}}\) is a free \(\mathbf{A}_{\mathfrak{p}}\) -module of rank \(n\) , \((\mathrm{P}_{(\mathbf{B})})_{\mathfrak{q}}\) is a free \(\mathbf{B}_{\mathfrak{q}}\) -module of rank \(n\) .

PROPOSITION 5. Let A be a semi-local ring and P a finitely generated projective A-module. If the rank of P is defined, P is a free A-module.

Suppose first that A is isomorphic to a product of fields \(K_{i}\) ( \(1 \leqslant i \leqslant n\) ). The \(K_{i}\) are then identified with the minimal ideals (Algebra, Chapter VIII, §3, no. 1) of A and, for all i, the sum \(p_{i}\) of the \(K_{j}\) of index \(j \neq i\) is a maximal ideal of A, the \(p_{i}\) ( \(1 \leqslant i \leqslant n\) ) being the only prime ideals of A. Every finitely generated A-module P is therefore the direct sum of its isotypical components \(P_{i}\) ( \(1 \leqslant i \leqslant n\) ), \(P_{i}\) being isomorphic to a direct sum of a finite number \(r_{i}\) of A-modules isomorphic to \(K_{i}\) (Algebra, Chapter VIII, §5, no. 1, Proposition 1 and no. 3, Proposition 11); the ring \(A_{p_{i}}\) is identified with \(K_{i}\) and annihilates the \(P_{j}\) of index \(j \neq i\) , hence \(r_{i} = \operatorname{rg}_{\mathfrak{p}_{i}}(\mathbf{P})\) ; if all the \(r_{i}\) are equal to the same number r, P is isomorphic to \(A^{r}\) , whence the proposition in this case. In the general case, let R be the Jacobson radical of A and B = A/R; as B is a product of fields, the projective B-module \(P_{(B)}\) is free by the remarks preceding Proposition 4. Also P is a flat A-module and the proposition then follows from §3, no. 2, Proposition 5.

